\CatchFileBetweenTags{\AlphaInvVal}{constants.tex}{AlphaInvVal}
\CatchFileBetweenTags{\AlphaInvMAR}{constants.tex}{AlphaInvMARVal}
\CatchFileBetweenTags{\MeMeVPrint}{constants.tex}{MeMeVPrint}

% CAP
\CatchFileBetweenTags{\HiggsVEVStepOneVal}{constants.tex}{HiggsVEVStepOneVal}
\CatchFileBetweenTags{\HiggsVEVStepTwoVal}{constants.tex}{HiggsVEVStepTwoVal}
\CatchFileBetweenTags{\HiggsVEVVal}{constants.tex}{HiggsVEVVal}
\CatchFileBetweenTags{\HiggsVEVExperimentalVal}{constants.tex}{HiggsVEVExperimentalVal}
\CatchFileBetweenTags{\HiggsVEVSigma}{constants.tex}{HiggsVEVSigma}
\CatchFileBetweenTags{\HiggsImpedanceVal}{constants.tex}{HiggsImpedanceVal}
\CatchFileBetweenTags{\WeakApertureProjVal}{constants.tex}{WeakApertureProjVal}

\CatchFileBetweenTags{\FermiConstVal}{constants.tex}{FermiConstVal}
\CatchFileBetweenTags{\FermiConstExperimentalVal}{constants.tex}{FermiConstExperimentalVal}
\CatchFileBetweenTags{\FermiConstSigma}{constants.tex}{FermiConstSigma}

% PRO
\CatchFileBetweenTags{\HiggsLambdaVal}{constants.tex}{HiggsLambdaVal}
\CatchFileBetweenTags{\HiggsLambdaStepOneVal}{constants.tex}{HiggsLambdaStepOneVal}
\CatchFileBetweenTags{\HiggsLambdaExperimentalVal}{constants.tex}{HiggsLambdaExperimentalVal}
\CatchFileBetweenTags{\HiggsLambdaSigma}{constants.tex}{HiggsLambdaSigma}

% TOL
\CatchFileBetweenTags{\HiggsMassVal}{constants.tex}{HiggsMassVal}
\CatchFileBetweenTags{\HiggsMassExperimentalVal}{constants.tex}{HiggsMassExperimentalVal}
\CatchFileBetweenTags{\HiggsMassSigma}{constants.tex}{HiggsMassSigma}

% MAR
\CatchFileBetweenTags{\ElectronYukawaVal}{constants.tex}{ElectronYukawaVal}
\CatchFileBetweenTags{\ElectronYukawaStepOneVal}{constants.tex}{ElectronYukawaStepOneVal}
\CatchFileBetweenTags{\ElectronYukawaExperimentalVal}{constants.tex}{ElectronYukawaExperimentalVal}
\CatchFileBetweenTags{\ElectronYukawaSigma}{constants.tex}{ElectronYukawaSigma}

\section{System 5: The Resonant Interface (The Electroweak Scale)}
\label{sec:resonant_interface}

The vacuum possesses a strictly finite channel capacity ($\nu=16$) operating at a fundamental temporal resonance ($\Delta=43$). In physical terms, this finite bandwidth imposes a strict ceiling on the local energy density of the lattice ($E = \hbar\omega$). The introduction of complex Topological Defects into the bulk saturates this core capacity, threatening the vacuum with Causal Aliasing: the routing failure where the substrate can no longer distinguish which channel belongs to which process when the capacity budget $\nu$ is exceeded, which manifests physically as runaway energy divergence.

To persist, the system must resolve this complexity. While System 4 defined the \textit{fractional ratio} of bandwidth allocated to this sector (e.g., $\sin^2\theta_W$), System 5 derives the \textit{absolute magnitude} of that capacity. This absolute ceiling manifests as the Electroweak Scale.

Under \textbf{Rule 3: Impedance Matching}, the vacuum must translate information between the deep geometric bulk (the core lattice) and the localized topological boundary (the particle surface) without dissipative reflection loss. To achieve this, the active energy bandwidth of the Weak Interaction must be perfectly matched to the bulk lattice resonance. Here, we derive the exact phenomenological constants of this interface as the partition coefficients of this geometric matching.

Consistent with \textit{Persistence Mechanics: Entropic Action} \cite{meyer_ieaction_2026}, we identify the Higgs sector as the mathematical bounding of this state-space capacity. The internal bandwidth of this resonant interface is rigidly partitioned across the six pillars, defining the following physical parameters:

\begin{itemize}
    \item \textbf{Capacity ($v$):} The VEV; the mass baseline allocated to the resonant interface, against which the vacuum potential is normalized.
    \item \textbf{Map ($Y$):} Scalar Hypercharge; the unitary address. It corresponds to the inverse of the lattice Unitary Ground State ($\Delta^0=1$).
    \item \textbf{Protocol ($\lambda$):} Self-Coupling; the self-interaction bandwidth allocation.
    \item \textbf{Governor ($\lambda|\phi|^4$):} The Quartic Potential; The geometric capacity wall enforcing the topological boundary ($\chi=2$).
    \item \textbf{Toll ($\mu^2$):} Instability Factor; the fractional entropic cost of maintaining the broken symmetry state.
    \item \textbf{Margin ($y_e$):} The Electron Yukawa; the minimal resolvable noise floor.
\end{itemize}

In the Standard Model, the seemingly arbitrary values of this sector constitute the ``Hierarchy Problem,'' but they are recognized here as the unique set of partition coefficients required to maintain vacuum stability against causal aliasing.

\subsection{Capacity: The Higgs VEV (\texorpdfstring{$v$}{v})}
Under \textbf{Rule 1: Capacity Partitioning}, the VEV represents the net informational capacity allocated to the Weak Interaction. We derive this by calculating the total available states over the causal boundary and subtracting the structural overhead.

\subsubsection{The Bare Geometric Floor}
To derive the absolute capacity of the Electroweak interface, we must first define the total informational envelope required to sustain a stable topological defect. Under \textbf{Rule 1: Capacity Partitioning}, the system must allocate enough bandwidth to account for every state transition within the defect's causal influence.

The temporal period $\Delta=43$ defines the \textit{causal horizon} of a lattice interaction (the distance information can travel before the system resets). Because information propagates isotropically at the speed of light ($c=1$), this resonance defines a 2-dimensional causal area $\Delta \times \Delta = \Delta^2$. (\textit{Scope Defense: The Euler characteristic $\chi=2$ uniquely identifies the defect boundary as a closed orientable 2-sphere ($S^2$) --- the only compact surface with $\chi=2$. $S^2$ is intrinsically 2-dimensional: its natural state-slot area is a 2D cross-section, not a $D$-dimensional volume. With $c=1$, both the radial and temporal horizons extend to $\Delta$, giving area $\Delta^2$. Using $\Delta^D = \Delta^4$ would count the full causal diamond volume rather than the 2-surface boundary --- an overcount by two dimensions. The choice of $\Delta^2$ is forced by the topology of the defect boundary, not assumed.})

To distinguish a specific particle (a defect) from the vacuum, this area must be bounded by the topological invariant $\chi=2$ (the sphere). The total capacity envelope required to resolve the defect's identity against the background is therefore the product of its topological signature and its resonant area: $\chi \Delta^2$. This represents the total number of ``state-slots'' available to the interface over one cycle.

The physical VEV is the energy scale of this interface capacity: the product of three substrate-level quantities --- (1) the available state-slots (the net capacity budget after infrastructure overhead), (2) the minimum energy per state ($m_e$, the Landauer limit derived below as the structural floor), and (3) the organizational resistance ($\alpha^{-1}$, the impedance governing how much energy is required to coherently address those states). From the total state-slot envelope, the substrate must reserve the infrastructure load it must process per cycle: ($L_{substrate} = \nu + D \cdot \Delta = 188$\cite{meyer_ieaction_2026}). Multiplying by the vacuum impedance ($\alpha^{-1}$) and the energy floor ($m_e$) yields the baseline energy scale before geometric corrections.

This transmission impedance is a strictly dimensionless ratio which can be mapped into a physical energy scale by anchoring it to the geometric resolution floor of the lattice: the minimum viable bit of mass ($m_e$, derived structurally below in Section \ref{subsec:margin_electron}). This establishes the true physical energy baseline of the field \textit{before} it is subjected to the macroscopic topological dilution and thermal noise corrections ($\delta_{topo}, \delta_{noise}$).

\begin{equation}
\begin{split}
v_{geo} &= (\chi \Delta^2 - L_{substrate}) \cdot \alpha^{-1} \cdot m_e \\
&= (2 \cdot 43^2 - 188) \cdot \AlphaInvVal \cdot \MeMeVPrint \text{ MeV} \\
&\approx \HiggsVEVStepOneVal
\end{split}
\end{equation}

\subsubsection{The Geometric Corrections}
The $v_{geo}$ baseline represents an idealized count of state-slots. To find the \textit{physical} energy scale, we must account for the reality of the projection: how this energy is distributed in continuous space and how it is resolved against the substrate noise floor. These effects manifest as exact leading-order linear expansions ($1 \pm x$), representing the first-order response of the vacuum to its own geometric constraints.

\paragraph{Topological Dilution (Spatial Correction)}
In the continuum limit, the localized charge of a defect is not a discrete point but a field distributed over a spherical gauge geometry ($4\pi$). This ``dilutes'' the available energy scale because the topological boundary ($\chi=2$) effectively expands the volume the coupling must cover. We define the \textbf{Effective Coupling Volume} ($D_{eff}$) as the manifold rank ($D=4$) augmented by this boundary-to-volume ratio:
\begin{equation}
D_{eff} = D + \frac{\chi}{4\pi} \approx 4.15915
\end{equation}
\textit{(Scope Defense: This continuous $4\pi$ correction represents the macroscopic solid angle of the $\chi=2$ boundary; it does not alter the integer dimensionality of the substrate. The correction is $\alpha/D_{eff}$ rather than $\alpha/D$ because the effective volume experienced by the coupling is $D_{eff}$, not the bare lattice rank $D$.)} This yields the dilution factor:
\begin{equation}
\delta_{topo} = \left( 1 + \frac{\alpha}{D_{eff}} \right)
\end{equation}

\paragraph{Thermodynamic Noise (Temporal Correction)}
We must also subtract the energy ``lost'' to the substrate's resolution floor. Under \textbf{Rule 2: Reversible Encoding}, the \textbf{Margin Impedance} ($Z_{MAR} \approx \AlphaInvMAR$) defines the minimum resolvable energy unit. Because the vacuum supports exactly 3 fermion generations (the available residual degrees of freedom: $\sigma - \chi = 3$), this noise is partitioned equally across these three independent ``signal slots.'' \textit{(Scope Defense: The divisor $(\sigma-\chi) = 3$ is the internal interaction bulk minus the topological boundary — the number of geometrically independent generation channels. Dividing by $\nu=16$ would conflate chiral capacity with generational noise channels, which are orthogonal degrees of freedom.)} To maintain global solvency, the physical VEV must be reduced by this noise amplitude:
\begin{equation}
\delta_{noise} = \left( 1 - \frac{Z_{MAR}}{\sigma - \chi} \right)
\end{equation}

\subsubsection{Synthesis: The Physical VEV}
The physical Vacuum Expectation Value is the geometric floor modulated by the simultaneous application of these topological and thermodynamic factors.
\begin{equation}
v_{phys} = v_{geo} \cdot \delta_{topo} \cdot \delta_{noise}
\end{equation}

\paragraph{What the Fermi window says about the corrections.} The measured Fermi constant resolves both factors. Without $\delta_{noise}$ the VEV lies above $246.21987$~GeV, outside the window (\texttt{VevCorrections.noise\_needed}); with divisor $k$ in place of $\sigma-\chi$ it lies inside exactly for $k=3$ or $k=4$ (\texttt{divisor\_window}). The topological factor has exactly the stated form: $e^{x}$ or $1/(1-x)$ in place of $1+x$, $x=\alpha/D_{eff}$, put $v$ above $246.22001$ and $246.22038$~GeV, $\alpha/D$ without the boundary term above $246.2366$~GeV, and the QED-like $3\alpha/(4\pi)$ below $246.2167$~GeV (\texttt{exp\_excluded}, \texttt{geometric\_excluded}, \texttt{bare\_dimension\_excluded}, \texttt{qed\_like\_excluded}). The window pins the effective dimension to $(4.1584,4.1599)$, and among boundary terms $a/(b\pi)$ with $b\le200$ only $a/b=1/2$, that is $\chi/(4\pi)$, fits (\texttt{dimension\_window}, \texttt{boundary\_unique}). Two persistence readings account for what the data select. The linear form is one repair: $1/(1-x)=\sum x^k$ is repairs of repairs and $e^x$ continuous compounding, while the anchor's maintenance is a single repair (\texttt{PersistenceClosures.one\_repair}), exactly first order. The divisor $3$ rather than $4$, which the data allow equally, is fixed by exact repair. The reference sits on the time axis; an exact repair leaves no residual on its axis, the residual is symmetric under the axis permutations that fix it, and so the noise is carried by the three spatial axes, $Y/3$ each (\texttt{NoiseShare.share\_spatial}, \texttt{symmetric\_of\_invariant}). Equal shares on all four axes would leave $Y/4$ on the reference's axis, contradicting exactness (\texttt{four\_needs\_residual}). The three are also the colour axes $\sigma-\chi$, the same fact. Both remaining readings now follow from exact repair and a normalized readout (\texttt{VevReadout}). An exact repair's net correction is the displacement itself, with no higher-order terms (\texttt{exact\_net\_is\_identity}), so one exact repair of the relative displacement $x$ gives $1+x$. Re-repairing a residual fraction $r$ would give $1+x/(1-r)$, which is $1+x$ only for the exact $r=0$ and is the excluded geometric form for $r=x$ (\texttt{residual\_series}). Replacing the one repair by $n\ge2$ equal partial repairs gives $(1+x/n)^n\ge1+x+x^2/4$, outside the Fermi window for every $n\ge2$ (\texttt{compound\_lower}, \texttt{compound\_excluded}), with $e^x$ the limit. For the noise, exact repair fixes the shares ($0$ on the reference's time axis, $Y/3$ on each spatial axis), so \emph{any} normalized readout law on the spatial axes reads $Y/3$ (\texttt{normalized\_share}); the VEV need not single out an axis. An unnormalized total would read $Y$, divisor $1$, which the window excludes (\texttt{total\_excluded}). A normalized law that also weights the time axis by $w_0$ reads $(1-w_0)Y/3$, and the uniform law over all four axes reads $Y/4$, which the window also allows (\texttt{time\_weight}, \texttt{uniform\_four}). The time axis drops out because a readout is an event. A comparison of the reference $\phi$ with its transport $U\phi$ fires with probability $\|\phi-U\phi\|^2/4$, which vanishes exactly when $U$ fixes $\phi$ (\texttt{VevEvent.flag\_zero\_iff}). Along the time axis the transport of the maintained reference is the maintenance update, which an exact repair returns to the reference, so that comparison never fires (\texttt{time\_never\_fires}). Conditioning any prior law over the four axes on a readout event therefore removes the time axis, and the event-conditioned reading is $Y/3$ (\texttt{conditioned\_on\_event}, \texttt{vev\_reads\_spatial\_share}). Both correction factors of the VEV now rest on exact repair, one repair and the event principle that nothing is registered without an event, the same principle used for depth (\texttt{DepthEvents.no\_event\_no\_advance}).

\paragraph{Which accounting the VEV can persist with.} The measured Fermi constant fixes $246.21957<v_F<246.21971$~GeV, finer than any accounting choice (\texttt{VevAccounting.fermi\_window}). The formula lies inside this window only with $L=D\Delta+\nu=188$, the 172 addresses and 16 channels counted as frame at unit weight (\texttt{frame\_inside}). Counting the 16 as stored payload, $172+16\eta$, puts $v$ above $246.2261$; subtracting the whole footprint $N_{sys}$ puts it below $245.876$ (\texttt{payload\_outside}, \texttt{footprint\_outside}). Any boundary share $s\ge10^{-3}$ of the wires is also excluded (\texttt{VevPremises.wire\_share\_excluded}): a wire is a 1-cell between nodes and meets a node's closed boundary only at points, so it occupies no boundary area. Under the event calibration the handoff factor is exactly $\alpha^{-1}=4\pi N_{sys}(\mathrm{toll}\cdot m)$, one footprint-wide repair over the full solid angle per interface slot (\texttt{handoff\_is\_footprint\_repair}). The premises that remain are that the electron's maintenance is one repair and the two correction factors $\delta_{topo}$ and $\delta_{noise}$, which the Fermi constant resolves individually. Why the electron supplies the anchor is no longer a separate choice. The shared electroweak reference is the lepton doublet (\cref{app:gauge_admittance}). On it the conserved charge $Q=T_3+Y/2$ takes the values $0$ and $-1$, so exactly one of its carriers is read by the unbroken charge, and it is the electron (\texttt{ReferenceAnchor.reference\_charges}, \texttt{electron\_unique\_charged\_reference}). The scale is anchored on the reference's electrically readable carrier. Its maintenance count is one repair: an unmaintained carrier does not persist, so the count is positive, and Economy takes the least positive count (\texttt{PersistenceClosures.one\_repair}). This fixes the count, not the physical content of the repair.

Substituting the derived values:
\begin{equation}
\begin{split}
v_{phys} &= (\HiggsVEVStepOneVal) \cdot \\
&\left( 1 + \frac{\alpha}{4.159} \right) \cdot \left( 1 - \frac{\AlphaInvMAR}{3} \right) \\
& \approx \mathbf{\HiggsVEVVal}
\end{split}
\end{equation}

\vspace{0.5em}
\begin{tabular}{ll}
\textbf{Prediction:} & $\HiggsVEVVal$ \\
\textbf{Experiment:} & \HiggsVEVExperimentalVal \\
\textbf{Accuracy:}   & $\HiggsVEVSigma \sigma$ \\
\end{tabular}
\vspace{0.5em}

\subsubsection{Output: The Fermi Constant (\texorpdfstring{$G_F$}{GF})}
With the interface capacity $v_{phys}$ fixed, a propagating fermion couples to it through a 1D chiral projection. The Fermi Constant $G_F$ is the inverse squared capacity of the interface, attenuated by this projection factor.

\paragraph{The Geometric Origin of $\sqrt{2}$}
For $v_{phys}$ to mediate a force, it must be projected onto the 1D chiral axis of the fermion’s trajectory. The interface amplitude is isotropically distributed on $S^1$; projecting onto a single axis averages as $\langle\cos^2\theta\rangle = 1/2$, yielding an RMS projection factor of exactly $1/\sqrt{2}$.

\paragraph{The Result}
\begin{equation}
G_F=\frac{1}{\sqrt{2}\,v_{phys}^2}\approx\FermiConstVal
\end{equation}
The same factor follows from the reference without the projection picture. On the reference's neutral carrier $\langle T_1^2\rangle=1/4$, so $m_W^2=g^2v^2/4$; exchanging the $W$ at low momentum gives the contact coefficient $g^2/(8m_W^2)$ (Paper~II, \texttt{PMAction.ContactMatching.contact\_limit}), and $G_F/\sqrt2=g^2/(8m_W^2)$ is then $1/(2v^2)$ for every $g$ (\texttt{PMNested.MatterRunning.fermi\_from\_reference}).

\vspace{0.5em}
\begin{tabular}{ll}
\textbf{Prediction:} & $\FermiConstVal$ \\
\textbf{Experiment:} & \FermiConstExperimentalVal \\
\textbf{Accuracy:}   & $\FermiConstSigma \sigma$ \\
\end{tabular}
\vspace{0.5em}

\subsubsection{Interface Matching and the Scalar Gain}
The values $Z_H\approx5.9644494$ and $A_{weak}\approx5.9651163$ are close but unequal. Lean certifies $1/2000<A_{weak}-Z_H<1/1000$. Under the stated positive-real interface law, the reflected fraction is
\begin{equation}
 10^{-9}<\left(\frac{A_{weak}-Z_H}{A_{weak}+Z_H}\right)^2<10^{-8}.
\end{equation}
Thus the same law has transmission greater than $0.99999999$ (\texttt{SurfaceBudget}). A reserve must actually be available to handle the reflected load; its small size is not itself a proof of solvency. The fraction is smaller than the proposed Margin coefficient, but allocating that coefficient to this interface would be an additional premise whose cost must appear elsewhere in the budget.

Maximal gain exhausts the stated 23-label regulatory ledger, leaving no reserve there. Exact aperture matching instead has a unique feasible gain in $(0.1294121,0.1294122)$, below $128/989$ (\texttt{MatchedRegulation}). Asking which gain persists separates the two. Under Closure the reflected load must be paid from the same ledger at some charge rate $c\ge0$: persistence at gain $g$ requires $cR(g)\le 3-(23g+1/43)$. For every $c>0$ maximal gain fails, because its reflection is positive and its slack is zero; exact matching persists but leaves slack unused. There is a unique largest persistent gain $g^\ast$, strictly between the matched gain and $128/989$, with
\begin{equation}
 \frac{128}{989}-g^\ast<\frac{c\cdot10^{-8}}{23}
\end{equation}
(\texttt{PersistentGain.persistent\_maximum}). Economy within persistence therefore selects maximal gain reduced by the reflection it must service; exact matching would be selected only for unbounded reflection cost. The charge rate follows from the readout that fixes the budget. A reflected amplitude $\sqrt R\,\psi$ has comparison response $R$ times that of $\psi$ (Paper~II, \texttt{RestorationResponse.response\_smul}; \texttt{ReflectionCharge.reflected\_response}), so restoring the reflected part of a complete repair costs $R$ times its readout, $3$, and $c=3$ (\texttt{charge\_rate}), given that the reflection is a fraction of the complete repair's displacement. Then $128/989-g^\ast<1.4\cdot10^{-9}$ (\texttt{gain\_at\_three}), and $\lambda$ and $m_H$ are unchanged at the displayed precision.

Neither side of this impedance comparison contains $v$, so the comparison cannot by itself select the electroweak mass scale. Nor does a numerical residual establish a universal minimum cost or a bound on omitted corrections. The formulas are retained as quantitative hypotheses with explicit joint-consistency tests.

\paragraph{A test of $\lambda$ without the VEV.} The reflection charge rate $c$ need not be derived: for every $0<c\le10^4$ the persistent gain lies within $4.4\times10^{-6}$ of $128/989$ (\texttt{ScalarSector.lambda\_insensitive}), far below the resolution of $\lambda$. With the measured Fermi constant, $v_F^2=1/(\sqrt2G_F)$ and $m_H=\sqrt{2\lambda}\,v_F$ lies in $(125.26,125.28)$~GeV (\texttt{higgs\_from\_fermi}), $0.5$--$0.75\sigma$ from $125.20\pm0.11$~GeV (\texttt{higgs\_pull}). This uses neither the VEV formula nor the electron mass. The ledger's $23=\nu+\sigma+\chi$ labels count the channels as frame; counting them as payload, $16\eta+7$, would give $m_H>125.52$~GeV, more than $2.8\sigma$ high (\texttt{ScalarAccounting.payload\_pull}).

\paragraph{The budget 3.} Among natural budgets only $B=3$ puts $(B-1/43)/23$ within $3\sigma$ of the measured $\lambda$ (\texttt{ScalarBudget.budget\_unique}). The selected electroweak reference supplies this count from the scalar's own task. Of the four real electroweak directions, exactly one (the multiples of $Q$) leaves the reference's neutral carrier fixed, and the other three displace it into a space of real dimension three (\texttt{ReferenceDirections.fixes\_iff}, \texttt{displacement\_dim}). Conservation requires the regulator to restore every direction in which the maintained reference can be displaced, and Economy excludes any in which it cannot, so one unit per moved direction gives $B=3$ (\texttt{budget\_is\_moved\_count}). That unit is the total gain of full restoration: a restoring map that returns every displacement within one update is the identity on the displacement space, whose trace is $3$ (\texttt{RegulatoryMaximin.full\_restore\_trace}), counted additively as the label side is. Equal gain on the labels is not a separate premise either. With $\sum_i g_i+1/43\le B$, some label receives at most $(B-1/43)/N$, and an allocation whose least gain reaches that value is the equal one (\texttt{min\_gain\_le}, \texttt{maximin\_forces\_equal}), so maximizing the worst label's margin gives equal gains. With $B=3$ the measured $\lambda$ then pins the count: only $N=23$ fits (\texttt{count\_unique}), although $\lambda$ alone does not fix the pair, since $(7,54)$ also fits (\texttt{pair\_not\_pinned}). Restoration need not complete within one update: any restoring map that returns every displacement after finitely many updates has $1-K$ nilpotent, so $\operatorname{tr}K=3$ (\texttt{ReferenceRestoration.finite\_restore\_trace}). Among sums of distinct architecture inputs $\{D,\Delta,\nu,N,\sigma,\chi,R_M,L_{\rm embed}\}$ the only one equal to $23$ is $\nu+\sigma+\chi$ (\texttt{unique\_decomposition}), so the measured coupling fixes which labels the ledger counts, given that they are distinct architecture inputs. That premise follows from what the scalar is. The scalar is the inter-sector link, and a link hop is indexed by a channel, an internal generator wire and a boundary unit, the apertures of the map pillar; its regulatory labels are these coordinate values, $\nu+\sigma+\chi=23$ (\texttt{RegulatoryLabels.label\_count}). To first order a hop's deviation is the sum of its coordinates' deviations, and regulating a label set removes every such displacement exactly when the set contains every label (\texttt{restores\_iff}); omitting one leaves the displacement concentrated on it (\texttt{missing\_label\_fails}). Regulating joint configurations instead would need $160$ labels, which Economy rejects since the coordinates suffice, and which the measured $\lambda$ excludes (\texttt{joint\_excluded}). Distinctness is itself the maximin choice: counting a label twice adds an entry and lowers every label's gain $(B-1/43)/N$ (\texttt{PersistenceClosures.repeat\_lowers\_margin}). Finally, the reference is maintained as an amplitude, not an occupation pattern: a diagonal reading, pure or mixed, commutes with both independent neutral generators $T_3$ and $Y/2$ (\texttt{diagonal\_reading\_keeps\_both}, \texttt{neutral\_independent}) and would leave two massless neutral bosons, whereas the amplitude reading leaves only $Q$. Exactness matters: a restorer that returns displacements only asymptotically, $K=k\cdot1$ with $0<k<2$, has total gain $3k$, any value in $(0,6)$ (\texttt{ReferenceLink.asymptotic\_trace\_free}). The measured $\lambda$ requires the exact restorer's $3$. Finiteness supplies part of this: on a finite state space a restorer under which every state converges to the reference reaches it exactly after a uniform number of updates (\texttt{PersistenceClosures.finite\_restorer\_exact}). Such a restorer cannot be a closed reversible update (\texttt{reversible\_cannot\_collapse}); it needs records or an environment. The link to the trace no longer needs a linear or nilpotent restorer. A complete repair, with any internal records and clock states, that returns every sufficiently small displacement exactly to the reference has net-correction derivative the identity, so its total restoring gain is $3$ (Paper~II, \texttt{PMAction.DifferentialRestoration.net\_gain\_trace}; \texttt{RepairBudget.repair\_gain}). Every displacement is that of an electroweak drift (\texttt{drift\_covers\_neighborhood}), so this exactness is the statement that every small drift is repaired; exactness on finitely many configurations would not fix the gain. Which functional of the repair the budget reads is not a further choice. Because the net correction of an exact repair is the identity on the displacement space (\texttt{BudgetReadout.repair\_derivative\_id}), any readout that adds over independent repairs, does not depend on the frame, and counts one restored direction as one unit gives $3$ (\texttt{budget\_readout}); the trace is one such readout (\texttt{trace\_is\_readout}), and counting the whole repair as a single event, $B=1$, is excluded by $\lambda$ (\texttt{single\_unit\_excluded}). The clock credit is a readout in the same unit: the clock phase is one direction, re-synchronized once per cycle, and a phase-invariant charge over the $43$-phase cycle is $1/43$ per phase (\texttt{ledger\_from\_readouts}). The coherent repair leaves its record occupied, and a conservative protocol that repairs $T$ drifts needs $T$ clean records (Paper~II, \texttt{RecordTransport.history\_packet\_bound}), so records must be renewed; but the record keeps at least the three restored directions, and charging its renewal to the budget per direction would put $\lambda$ outside its window (\texttt{renewal\_per\_direction\_excluded}). So the $1/43$ is not a per-direction record renewal. In Paper~II's walk implementation the record is a control channel of the walk carrying the doublet on its own payload index, on which weak acts faithfully (the walk's spinor coordinates alone admit only Abelian actions), and it is renewed by the incoming stream at the price of one seed preparation; charged once per cycle and spread over the phases, that price leaves $\lambda$ inside its window, so the data do not decide whether renewal is drawn from the regulatory budget. What supplies the incoming stream remains open. The colour-axis count $\sigma-\chi$ and the generation count are also $3$; the observable cannot separate them, but only the moved-direction count refers to what the scalar regulates.

\paragraph{How the link sets fermion masses.} Paper~II's detuned domain wall gives a light fermion with $m(m+1)=\delta^2$ in cutoff units. If the link fixes the detuning cost, $\delta^2=yv$, the Higgs coupling modifier $\kappa=(v/m)\,dm/dv$ is $(m+1)/(2m+1)$, within $(0.99,1]$ for fermions far below the cutoff; if it fixes the amplitude, $\delta=yv$, then $\kappa=2(m+1)/(2m+1)\approx2$ (\texttt{FermionMassLaw.law\_cost}, \texttt{law\_amp}). The measured $\kappa_\tau=0.92\pm0.08$ places the first law within $1\sigma$ and the second more than $13\sigma$ away (\texttt{tau\_decides}). The link therefore sets the detuning cost, consistent with the event law's pricing of squared amplitudes, and the mass law is $m(m+1)=yv$: the linear law with an $O(m)$ correction.

\CatchFileBetweenTags{\ElectronClosureCoefficientVal}{constants.tex}{ElectronClosureCoefficientVal}
\CatchFileBetweenTags{\ElectronYukawaSigmaDVal}{constants.tex}{ElectronYukawaSigmaDVal}

\subsubsection{Margin: The Electron Yukawa}
\label{subsec:margin_electron}
With the VEV fixed, the electron Yukawa follows from the mass law $m_e=y_e v/\sqrt{2}$:
\begin{equation}
 y_e=\frac{\sqrt{2}\,m_e}{v_{phys}}\approx\ElectronYukawaVal,
\end{equation}
so the electron relation closes exactly (Lean: \texttt{ScaleConsistency.closure\_exact}). This $y_e$ is not an independent prediction: the dimensionless content is the VEV ratio $v/m_e$, which reproduces the Fermi-constant value of $v$ well within its uncertainty.

\vspace{0.5em}
\begin{tabular}{ll}
\textbf{Mass law from $v$:} & $\ElectronYukawaVal$ \\
\textbf{Tree-level $\sqrt2 m_e/v$:} & \ElectronYukawaExperimentalVal \\
\end{tabular}
\vspace{0.5em}

An independent dressing of the Margin floor, $y_e=PM(1+c\,\alpha)$ with $PM\approx\ElectronYukawaStepOneVal$, is consistent with the VEV only for $1.28777<c<1.28779$; the calculator gives $c=\ElectronClosureCoefficientVal$ (\texttt{closure\_coefficient}). The earlier proposal $c=\sigma/D=5/4$ gives $\ElectronYukawaSigmaDVal$ and is excluded by this joint check (\texttt{stated\_dressing\_excluded}); a scalar-field rescaling or a $\sqrt2$ convention cannot repair it (\texttt{RepairSpace}). A derived coefficient near $1.2878$ would restore an independent Margin prediction; no such derivation is claimed.
